\documentclass[12pt]{article}
\usepackage[utf8]{inputenc}
\usepackage[T1]{fontenc}
\usepackage[spanish]{babel}
\usepackage{amsmath, amssymb, amsfonts}
\usepackage{fancyhdr}
\usepackage{xcolor}
\usepackage[margin=3cm]{geometry}
\usepackage{graphicx}
\usepackage{float}
\usepackage{multicol}
\usepackage{enumitem}
\usepackage{tcolorbox}

% Encabezado
\pagestyle{fancy}
\fancyhf{}
\fancyhead[L]{Escuela de Ingeniería \\ Universidad de O'Higgins}
\fancyhead[R]{Analisis y Diseño de Software \\ 2026-01}
\fancyfoot[C]{\thepage}

% Título comprimido
\makeatletter
\renewcommand{\maketitle}{%
  \begin{center}
    {\LARGE \@title\par}
    \vskip 0.3em
    {\large \@author\par}
    \vskip 0.3em
  \end{center}
  \vskip -0.2em  % Espacio después del título
}
\makeatother

\title{Ayudantía 5}
\author{Felipe Llach R.}

\begin{document}
\maketitle
\section{Contexto del Proyecto}
Una emergente cadena de \textit{food trucks} ha detectado un cuello de botella en su operación diaria: la comunicación entre la toma de pedidos y la preparación en cocina es lenta y propensa a errores debido al uso de tickets de papel.

Para solucionar esto, el equipo de desarrollo debe construir un \textbf{Kitchen Display System (KDS)}. El flujo esperado de la información es el siguiente:

\begin{enumerate}
    \item El cajero toma el pedido utilizando el Sistema de Punto de Venta (POS) existente de la empresa.
    \item El POS envía la información del pedido al nuevo sistema KDS.
    \item Los cocineros visualizan los pedidos en tiempo real a través de tablets instaladas en cada \textit{food truck}.
    \item Desde la tablet, los cocineros pueden marcar los platos como ``En preparación'' y ``Completados''.
    \item Un administrador central necesita poder extraer métricas de tiempos de preparación al final del día.
\end{enumerate}

Su tarea como arquitectos de software es modelar esta solución utilizando los primeros tres niveles del Modelo C4.

\section{Instrucciones de Desarrollo}
A continuación, resuelva el ejercicio diseñando los diagramas correspondientes.
\subsection{Parte 1: Diagrama de Contexto (Nivel 1)}
\textbf{Objetivo:} Modelar el panorama general del sistema y sus interacciones con el mundo exterior.
\begin{itemize}
    \item Identifique al sistema principal (KDS).
    \item Identifique a todos los actores (usuarios humanos) que interactuarán directa o indirectamente con el KDS.
    \item Identifique cualquier sistema de software externo involucrado en el flujo de trabajo.
    \item \textbf{Entregable:} Un diagrama de Contexto que muestre el sistema KDS en el centro, rodeado de sus usuarios y dependencias externas, con flechas direccionales que expliquen brevemente qué información fluye entre ellos.
\end{itemize}

\subsection{Parte 2: Diagrama de Contenedores (Nivel 2)}
\textbf{Objetivo:} Definir la arquitectura de alto nivel del sistema KDS, mostrando sus aplicaciones y unidades de almacenamiento de datos.
\begin{itemize}
    \item Haga un ``zoom'' dentro del sistema KDS.
    \item Diseñe los contenedores necesarios para que el sistema funcione. Considere que el equipo de desarrollo tiene experiencia trabajando con \textbf{C\# para el backend}.
    \item Debe incluir, como mínimo, la interfaz de usuario de los cocineros, la lógica de negocio central (API) y un mecanismo de persistencia de datos (Base de Datos).
    \item Indique las tecnologías sugeridas para cada contenedor (ej. tipo de base de datos, framework del backend, plataforma del frontend).
    \item \textbf{Entregable:} Un diagrama de Contenedores detallando las piezas desplegables del KDS, las responsabilidades de cada contenedor y los protocolos de comunicación entre ellos.
\end{itemize}

\subsection{Parte 3: Diagrama de Componentes (Nivel 3)}
\textbf{Objetivo:} Estructurar el código internamente, aplicando principios de diseño orientado a objetos.
\begin{itemize}
    \item Haga un ``zoom'' \textbf{exclusivamente dentro de la API desarrollada en C\#}.
    \item Modele los componentes internos necesarios para recibir un nuevo pedido desde el POS y para actualizar el estado del pedido desde la tablet del cocinero.
    \item Considere separar las responsabilidades utilizando modelos, controladores y repositorios.
    \item \textbf{Entregable:} Un diagrama de Componentes del contenedor de la API, mostrando cómo interactúan los controladores, servicios y repositorios entre sí, y cómo se conectan hacia afuera con la base de datos y los otros contenedores.
\end{itemize}

\end{document}